%%% !TeX document-id = {7f856805-4d41-4d13-829b-b08b1a8ea49a}
 %%%  کلاس AUTthesis، نسخه اردیبهشت 1396
%%%   دانشگاه صنعتی امیرکبیر                 http://www.aut.ac.ir
%%%  تالار گفتگوی پارسی‌لاتک،       http://forum.parsilatex.com
%%%   آپدیت شده در آبان 95
%%      پشتیبانی و راهنمایی          badali_farhad@yahoo.com
%%%
%%%   بازبینی و اصلاح شده در اردیبهشت ماه 1396
%%%  Tested via TeXstudio in TeXlive 2014,‌‌ 2015 & 2016.
%%%

%-----------------------------------------------------------------------------------------------------
%        روش اجرا.: 2 بار F1 ، 2 بار  F11(به منظور تولید مراجع) ، دوبار Ctrl+Alt+I (به منظور تولید نمایه) و دو بار F1 -------> مشاهده Pdf
%%%%%%%%%%%%%%%%%%%%%%%%%%%%%%%%%%%%%%%%%%%%%%%%%%%%%%
%   TeXstudio as your IDE
%%  برای compile در TeXstudio تنها کافی است منوی Options->Configure TeXstudio را زده و در پنجره Configure TeXstudio در بخش Build گزینه Default Compiler را به XeLaTeX تغییر دهید. سند شما به راحتی compile خواهد شد.
%   F1 & F5 : Build & view
%   F6      : Compile
%   F7      : View
%   --------------
%   شما می‌توانید این نرم‌افزار را از سایت زیر دریافت کنید. برای استفاده از این نرم‌افزار باید پیش از نصب آن TeXlive را نصب کرده باشید.
%   TeXstudio.org
%   این سند با نسخه‌های 2014، 2015 و 2016 نرم‌افزار TeXlive در محیط TeXstudio آزمایش شده است.
%   نکته: در هنگام compile با نرم‌افزار TeXstudio فایل اصلی حتما باید باز باشد ولی نیازی نیست که روی tab آن باشید.
%	همچنین در پوشه LookAtMe فایلهای png برای راهنمایی وجود دارند.
%	برای ساخت میانبر نیم فاصله نیز به فایلهای پوشه LookAtMe را نگاه کنید.
%   باآرزوی پیروزی برای شما
%   آریا
%%%%%%%%%%%%%%%%%%%%%%%%%%%%%%%%%%%%%%%%%%%%%%%%%%%%%%
%        اگر قصد نوشتن رساله دکتری را دارید، در خط زیر به جای msc،
%      کلمه phd را قرار دهید. کلیه تنظیمات لازم، به طور خودکار، اعمال می‌شود.
%%% !TEX TS-program = XeLaTeX
\documentclass[oneside,msc]{AUTthesis}
%       فایل commands.tex را حتماً به دقت مطالعه کنید؛ چون دستورات مربوط به فراخوانی بسته زی‌پرشین 
%       و دیگر بسته‌ها و ... در این فایل قرار دارد و بهتر است که با نحوه استفاده از آنها آشنا شوید. توجه شود برای نسخه نهایی پایان‌نامه حتماً hyperref را 
%        غیرفعال کنید.


% در این فایل، دستورها و تنظیمات مورد نیاز، آورده شده است.
%-------------------------------------------------------------------------------------------------------------------
% در ورژن جدید زی‌پرشین برای تایپ متن‌های ریاضی، این سه بسته، حتماً باید فراخوانی شود.
\usepackage{amsthm,amssymb,amsmath,amsfonts}
% بسته‌ای برای تنطیم حاشیه‌های بالا، پایین، چپ و راست صفحه
\usepackage[top=30mm, bottom=30mm, left=25mm, right=30mm]{geometry}
% بسته‌‌ای برای ظاهر شدن شکل‌ها و تصاویر متن
\usepackage{graphicx}
\usepackage{color}
%بسته‌ای برای تنظیم فاصله عمودی خط‌های متن
\usepackage{setspace}
\setstretch{1.3}
\usepackage{tocstyle}
\usetocstyle{standard}
\usepackage{titleps}
\usepackage{titletoc}
\usepackage{tocloft}
\usepackage{multirow}
\usepackage{tabularx}
\usepackage{booktabs}
%با فعال کردن بسته زیر فوت‌نوت‌ها در هر صفحه ریست می‌شوند. حالت پیش‌فرض آن ریست شدن در هر فصل می‌باشد.
%\usepackage[perpage]{footmisc}
\usepackage{enumitem}
%\usepackage{titlesec}
% بسته‌ و دستوراتی برای ایجاد لینک‌های رنگی با امکان جهش
\usepackage[pagebackref=false,colorlinks,linkcolor=blue,citecolor=red]{hyperref}
\usepackage[nameinlink]{cleveref}%capitalize,,noabbrev
 \AtBeginDocument{%
    \crefname{equation}{برابری}{equations}%
    \crefname{chapter}{فصل}{chapters}%
    \crefname{section}{بخش}{sections}%
    \crefname{appendix}{پیوست}{appendices}%
    \crefname{enumi}{مورد}{items}%
    \crefname{footnote}{زیرنویس}{footnotes}%
    \crefname{figure}{شکل}{figures}%
    \crefname{table}{جدول}{tables}%
    \crefname{theorem}{قضیه}{theorems}%
    \crefname{lemma}{لم}{lemmas}%
    \crefname{corollary}{نتیجه}{corollaries}%
    \crefname{proposition}{گزاره}{propositions}%
    \crefname{definition}{تعریف}{definitions}%
    \crefname{result}{نتیجه}{results}%
    \crefname{example}{مثال}{examples}%
    \crefname{remark}{نکته}{remarks}%
    \crefname{note}{یادداشت}{notes}%
}
% چنانچه قصد پرینت گرفتن نوشته خود را دارید، خط بالا را غیرفعال و  از دستور زیر استفاده کنید چون در صورت استفاده از دستور زیر‌‌، 
% لینک‌ها به رنگ سیاه ظاهر خواهند شد که برای پرینت گرفتن، مناسب‌تر است
%\usepackage[pagebackref=false]{hyperref}
% بسته‌ لازم برای تنظیم سربرگ‌ها
\usepackage{fancyhdr}
% بسته‌ای برای ظاهر شدن «مراجع»  در فهرست مطالب
\usepackage[nottoc]{tocbibind}
% دستورات مربوط به ایجاد نمایه
\usepackage{makeidx,multicol}
\setlength{\columnsep}{1.5cm}
%\usepackage{tensor}
\usepackage[all]{xy}
\usepackage{tikz}
\usetikzlibrary{arrows}
%%%%%%%%%%%%%%%%%%%%%%%%%%
\usepackage{verbatim}
\makeindex
\usepackage{sectsty}
% فراخوانی بسته زی‌پرشین و تعریف قلم فارسی و انگلیسی
\usepackage{xepersian}%[extrafootnotefeatures]
\SepMark{-}
%حتماً از تک لایو 2014 استفاده کنید.
\settextfont[Scale=1.1,Path=fonts/]{BNazanin}
\setlatintextfont{Times New Roman}
\renewcommand{\labelitemi}{$\bullet$}
%%%%%%%%%%%%%%%%%%%%%%%%%%
% چنانچه می‌خواهید اعداد در فرمول‌ها، انگلیسی باشد، خط زیر را غیرفعال کنید.
%در غیر اینصورت حتماً فونت PGaramond را نصب کنید.
%\setdigitfont[Scale=1.1]{PGaramond}%%Yas
%%%%%%%%%%%%%%%%%%%%%%%%%%
% تعریف قلم‌های فارسی اضافی برای استفاده در بعضی از قسمت‌های متن
\defpersianfont\nastaliq[Scale=2,Path=fonts/]{IranNastaliq}
\defpersianfont\chapternumber[Scale=3,Path=fonts/]{BNazanin}
\chapterfont{\centering}%
%%%%%%%%%%%%%%%%%%%%%%%%%%
% دستوری برای تغییر نام کلمه «اثبات» به «برهان»
\renewcommand\proofname{\textbf{برهان}}

% دستوری برای تغییر نام کلمه «کتاب‌نامه» به «منابع و مراجع«
\renewcommand{\bibname}{منابع و مراجع}


% Headings for every page of ToC, LoF and Lot
\setlength{\cftbeforetoctitleskip}{-1.2em}
\setlength{\cftbeforelottitleskip}{-1.2em}
\setlength{\cftbeforeloftitleskip}{-1.2em}
\setlength{\cftaftertoctitleskip}{-1em}
\setlength{\cftafterlottitleskip}{-1em}
\setlength{\cftafterloftitleskip}{-1em}
%%\makeatletter
%%%%\renewcommand{\l@chapter}{\@dottedtocline{1}{1em\bfseries}{1em}}
%%%%\renewcommand{\l@section}{\@dottedtocline{2}{2em}{2em}}
%%%%\renewcommand{\l@subsection}{\@dottedtocline{3}{3em}{3em}}
%%%%\renewcommand{\l@subsubsection}{\@dottedtocline{4}{4em}{4em}}
%%%%\makeatother


\newcommand\tocheading{\par عنوان\hfill صفحه \par}
\newcommand\lofheading{\hspace*{.5cm}\figurename\hfill صفحه \par}
\newcommand\lotheading{\hspace*{.5cm}\tablename\hfill صفحه \par}

\renewcommand{\cftchapleader}{\cftdotfill{\cftdotsep}}
\renewcommand{\cfttoctitlefont}{\hspace*{\fill}\LARGE\bfseries}%\Large
\renewcommand{\cftaftertoctitle}{\hspace*{\fill}}
\renewcommand{\cftlottitlefont}{\hspace*{\fill}\LARGE\bfseries}%\Large
\renewcommand{\cftafterlottitle}{\hspace*{\fill}}
\renewcommand{\cftloftitlefont}{\hspace*{\fill}\LARGE\bfseries}
\renewcommand{\cftafterloftitle}{\hspace*{\fill}}

%%%%%%%%%%%%%%%%%%%%%%%%%%
% تعریف و نحوه ظاهر شدن عنوان قضیه‌ها، تعریف‌ها، مثال‌ها و ...
%برای شماره گذاری سه تایی قضیه ها
\theoremstyle{definition}
\newtheorem{definition}{تعریف}[section]
\newtheorem{remark}[definition]{نکته}
\newtheorem{note}[definition]{یادداشت}
\newtheorem{example}[definition]{نمونه}
\newtheorem{question}[definition]{سوال}
\newtheorem{remember}[definition]{یاداوری}
\theoremstyle{theorem}
\newtheorem{theorem}[definition]{قضیه}
\newtheorem{lemma}[definition]{لم}
\newtheorem{proposition}[definition]{گزاره}
\newtheorem{corollary}[definition]{نتیجه}
%%%%%%%%%%%%%%%%%%%%%%%%
%%%%%%%%%%%%%%%%%%%
%%% برای شماره گذاری چهارتایی قضیه ها و ...
%%\newtheorem{definition1}[subsubsection]{تعریف}
%%\newtheorem{theorem1}[subsubsection]{قضیه}
%%\newtheorem{lemma1}[subsubsection]{لم}
%%\newtheorem{proposition1}[subsubsection]{گزاره}
%%\newtheorem{corollary1}[subsubsection]{نتیجه}
%%\newtheorem{remark1}[subsubsection]{نکته}
%%\newtheorem{example1}[subsubsection]{مثال}
%%\newtheorem{question1}[subsubsection]{سوال}

%%%%%%%%%%%%%%%%%%%%%%%%%%%%

% دستورهایی برای سفارشی کردن صفحات اول فصل‌ها
\makeatletter
\newcommand\mycustomraggedright{%
 \if@RTL\raggedleft%
 \else\raggedright%
 \fi}
\def\@makechapterhead#1{%
\thispagestyle{style1}
\vspace*{20\p@}%
{\parindent \z@ \mycustomraggedright %\@mycustomfont
\ifnum \c@secnumdepth >\m@ne
\if@mainmatter

\bfseries{\Huge \@chapapp}\small\space {\chapternumber\thechapter}
\par\nobreak
\vskip 0\p@
\fi
\fi
\interlinepenalty\@M 
\Huge \bfseries #1\par\nobreak
\vskip 120\p@

}

%\thispagestyle{empty}
\newpage}
\bidi@patchcmd{\@makechapterhead}{\thechapter}{\tartibi{chapter}}{}{}
\bidi@patchcmd{\chaptermark}{\thechapter}{\tartibi{chapter}}{}{}
\makeatother

\pagestyle{fancy}
\renewcommand{\chaptermark}[1]{\markboth{\chaptername~\tartibi{chapter}: #1}{}}

\fancypagestyle{style1}{
\fancyhf{} 
\fancyfoot[c]{\thepage}
\fancyhead[R]{\leftmark}%
\renewcommand{\headrulewidth}{1.2pt}
}


\fancypagestyle{style2}{
\fancyhf{}
\fancyhead[R]{چکیده}
\fancyfoot[C]{\thepage{}}
\renewcommand{\headrulewidth}{1.2pt}
}

\fancypagestyle{style3}{%
  \fancyhf{}%
  \fancyhead[R]{فهرست نمادها}
  \fancyfoot[C]{\thepage}%
  \renewcommand{\headrulewidth}{1.2pt}%
}

\fancypagestyle{style4}{%
  \fancyhf{}%
  \fancyhead[R]{فهرست جداول}
  \fancyfoot[C]{\thepage}%
  \renewcommand{\headrulewidth}{1.2pt}%
}

\fancypagestyle{style5}{%
  \fancyhf{}%
  \fancyhead[R]{فهرست اشکال}
  \fancyfoot[C]{\thepage}%
  \renewcommand{\headrulewidth}{1.2pt}%
}

\fancypagestyle{style6}{%
  \fancyhf{}%
  \fancyhead[R]{فهرست مطالب}
  \fancyfoot[C]{\thepage}%
  \renewcommand{\headrulewidth}{1.2pt}%
}

\fancypagestyle{style7}{%
  \fancyhf{}%
  \fancyhead[R]{نمایه}
  \fancyfoot[C]{\thepage}%
  \renewcommand{\headrulewidth}{1.2pt}%
}

\fancypagestyle{style8}{%
  \fancyhf{}%
  \fancyhead[R]{منابع و مراجع}
  \fancyfoot[C]{\thepage}%
  \renewcommand{\headrulewidth}{1.2pt}%
}
\fancypagestyle{style9}{%
  \fancyhf{}%
  \fancyhead[R]{واژه‌نامه‌ی فارسی به انگلیسی}
  \fancyfoot[C]{\thepage}%
  \renewcommand{\headrulewidth}{1.2pt}%
}
%


%دستور حذف نام لیست تصاویر و لیست جداول از فهرست مطالب
\newcommand*{\BeginNoToc}{%
  \addtocontents{toc}{%
    \edef\protect\SavedTocDepth{\protect\the\protect\value{tocdepth}}%
  }%
  \addtocontents{toc}{%
    \protect\setcounter{tocdepth}{-10}%
  }%
}
\newcommand*{\EndNoToc}{%
  \addtocontents{toc}{%
    \protect\setcounter{tocdepth}{\protect\SavedTocDepth}%
  }%
}
\newcounter{savepage}
\renewcommand{\listfigurename}{فهرست اشکال}
\renewcommand{\listtablename}{فهرست جداول}
%\renewcommand\cftsecleader{\cftdotfill{\cftdotsep}}
%%%%%%%%%%%%%%%%%%%%%%%%%%%%%
%%%%%%%%%%%%%%%%%%%%%%%%%%%%

\begin{document}
\baselineskip=.75cm
%% -!TEX root = AUTthesis.tex
% در این فایل، عنوان پایان‌نامه، مشخصات خود، متن تقدیمی‌، ستایش، سپاس‌گزاری و چکیده پایان‌نامه را به فارسی، وارد کنید.
% توجه داشته باشید که جدول حاوی مشخصات پروژه/پایان‌نامه/رساله و همچنین، مشخصات داخل آن، به طور خودکار، درج می‌شود.
%%%%%%%%%%%%%%%%%%%%%%%%%%%%%%%%%%%%
% دانشکده، آموزشکده و یا پژوهشکده  خود را وارد کنید
\faculty{دانشکده مهندسی کامپیوتر و فناوری اطلاعات}
% گرایش و گروه آموزشی خود را وارد کنید
\department{گرایش شبکه‌های کامپیوتری}
% عنوان پایان‌نامه را وارد کنید
\fatitle{زنجیره‌سازی کارکردهای مجازی سرویس شبکه با لحاظ محدودیت منابع مدیریتی}
% نام استاد(ان) راهنما را وارد کنید
\firstsupervisor{دکتر بهادر بخشی}
%\secondsupervisor{استاد راهنمای دوم}
% نام استاد(دان) مشاور را وارد کنید. چنانچه استاد مشاور ندارید، دستور پایین را غیرفعال کنید.
% \firstadvisor{نام کامل استاد مشاور}
%\secondadvisor{استاد مشاور دوم}
% نام نویسنده را وارد کنید
\name{پرهام}
% نام خانوادگی نویسنده را وارد کنید
\surname{الوانی}
%%%%%%%%%%%%%%%%%%%%%%%%%%%%%%%%%%
\thesisdate{شهریور ۱۳۹۸}

% چکیده پایان‌نامه را وارد کنید
\fa-abstract{
    مساله‌ی مجازی سازی توابع شبکه سعی دارد توابع شبکه را به صورت مجازی در شبکه جایگذاری نمایند و در ادامه
    با برقراری ارتباط میان آن‌ها سرویس‌هایی را فراهم آورد.
    یکی از مسائل در این روش پذیرش سرویس‌ها و قرار دادن آن‌ها بر روی زیرساخت است
    که در کارهایی زیادی به آن پرداخته شده است ولی یکی از اجزا معماری مجازی سازی کارکردهای
    شبکه بخش مدیریتی است که می‌بایست در کنار سرویس‌ها بر روی زیرساخت مستقر شود.
    در این رساله ما قصد داریم جایگذاری سرویس‌ها با لحاظ منابع مدیریتی را
    مدل‌سازی و حل نماییم.
}


% کلمات کلیدی پایان‌نامه را وارد کنید
\keywords{مجازی سازی کارکردهای شبکه، زنجیره‌سازی کارکردهای مجازی سرویس شبکه،بهینه‌سازی، بهینه‌سازی خطی صحیح}



\AUTtitle
%%%%%%%%%%%%%%%%%%%%%%%%%%%%%%%%%%
\vspace*{7cm}
\thispagestyle{empty}
\begin{center}
\includegraphics[height=5cm,width=12cm]{images/in-the-name-of-god}
\end{center}
% تاییدیه دفاع
\newpage
\thispagestyle{empty}
%\fontsize{18pt}{19pt}\selectfont

\section*{صفحه فرم ارزیابی و تصویب پایان نامه- فرم تأیید اعضاء كميته دفاع}

\fontsize{12pt}{14pt}\selectfont
\renewcommand{\baselinestretch}{1.5}
\vspace*{1cm}
   در این صفحه فرم دفاع یا تایید و تصویب پایان نامه موسوم به فرم کمیته دفاع- موجود در پرونده آموزشی- را قرار دهید.
\vspace*{1cm}


\subsection*{نکات مهم:}
 
\begin{itemize}
\item
	نگارش پایان نامه/رساله باید به
	{\color{red}
		زبان فارسی
	}
	و بر اساس آخرین نسخه دستورالعمل و راهنمای تدوین پایان نامه های دانشگاه صنعتی امیرکبیر باشد.(دستورالعمل و راهنمای حاضر)
\item رنگ جلد پایان نامه/رساله چاپي كارشناسي، كارشناسي ارشد و دكترا  بايد به ترتيب مشكي، طوسي و سفيد رنگ باشد.  
\item چاپ و صحافی پایان نامه/رساله بصورت
{\color{red}
	پشت و رو(دورو)
}
بلامانع است و انجام آن توصيه مي شود. 
\end{itemize}
%%%%%%%%%%%%%%%%%%%%%%%%%%%%%%%%%%%%%%%%%%%%%%%%%%%%%%%%%%%%%%%%%%%%%%%%%%%%%%%%%%%%%%%%%%%%%%%%%%
%%%%%%%%%%%%%%%%%%%%%%%%%%%%%%%%%%%%%%%%%%%%%%%%%%%%%%%%%%%%%%%%%%%%%%%%%%%%%%%%%%%%%%%%%%%%%%%%%%
\newpage
\thispagestyle{empty}
\begin{picture}(50,50)
  \put(10,0){\includegraphics[scale=.06]{images/fa-logo}}
  \put(4.5,-13){\footnotesize{دانشگاه صنعتی امیرکبیر}}
  \put(10.5,-27){\footnotesize{(پلی‌تکنیک تهران)}}
  \put(170,30){\bf{به نام خدا}}
  \put(140,-5){\Large\bf{تعهدنامه اصالت اثر}}
  \put(300,0){تاریخ: \datethesis}
\end{picture}

\vspace*{2.5cm}

اينجانب {\bf{\fname\ \lname}} متعهد می‌شوم که مطالب مندرج در این پایان‌نامه حاصل کار پژوهشی اینجانب تحت نظارت و راهنمایی اساتید دانشگاه صنعتی امیرکبیر بوده و به دستاوردهای دیگران که در این پژوهش از آنها استفاده شده است مطابق مقررات و روال متعارف ارجاع و در فهرست منابع و مآخذ ذکر گردیده است. این پایان‌نامه قبلاً برای احراز هیچ مدرک هم‌سطح یا بالاتر ارائه نگردیده است.

در صورت اثبات تخلف در هر زمان، مدرک تحصیلی صادر شده توسط دانشگاه از درجه اعتبار ساقط بوده و دانشگاه حق پیگیری قانونی خواهد داشت.


کلیه نتایج و حقوق حاصل از این پایان‌نامه متعلق به دانشگاه صنعتی امیرکبیر می‌باشد. هرگونه استفاده از نتایج علمی و عملی، واگذاری اطلاعات به دیگران یا چاپ و تکثیر، نسخه‌برداری، ترجمه و اقتباس از این پایان نامه بدون موافقت کتبی دانشگاه صنعتی امیرکبیر ممنوع است. 
نقل مطالب با ذکر مآخذ بلامانع است.\\
\vspace{2.5cm}


{\centerline {\bf{\fname\ \lname}}}
\vspace*{.2cm}
{\centerline{امضا}}
%%%%%%%%%%%%%%%%%%%%%%%%%%%%%%%%%
% چنانچه مایل به چاپ صفحات «تقدیم»، «نیایش» و «سپاس‌گزاری» در خروجی نیستید، خط‌های زیر را با گذاشتن ٪  در ابتدای آنها غیرفعال کنید.
% پایان‌نامه خود را تقدیم کنید
% نیایش خود را در فایل زیر بنویسید.
\begin{acknowledgementpage}

\vspace{1.5cm}

{\nastaliq
{
 نويسنده پايان‌نامه، درصورت تمايل ميتواند برای سپاسگزاری پايان‌نامه خود را به شخص يا اشخاص و يا ارگان خاصی تقدیم نماید.
}}\end{acknowledgementpage}
\newpage
%سپاسگزاری را در فایل زیر بنویسید.
%%%%%%%%%%%%%%%%%%%%%%%%%%%%%%%%%%%%
\newpage\thispagestyle{empty}
% سپاس‌گزاری
{\nastaliq
سپاس‌گزاری
}
\\[2cm]
در اینجا لازم می‌دانم از راهنمایی‌ها و مساعدت‌های اساتید عزیز و گرانقدرم
جناب آقای دکتر بخشی صمیمانه قدردانی و سپاس‌گزاری نمایم.
در ادامه از دوست خوبم بهروز فرکیانی که همواره من را راهنمایی کرده‌
و
از پدر و مادرم که همواره من را حمایت کرده‌اند تشکر می‌کنم.


در نهایت جا دارد از دوست، همکار و مدیر خوبم سینا سعیدی مدیریت فنی تیم منابع مشترک شرکت ایده گزین ارتباطات روماک تشکر کنم که بدون حمایت‌های ایشان نگارش این پایان نامه ممکن نبود.








% با استفاده از دستور زیر، امضای شما، به طور خودکار، درج می‌شود.
\signature








%%%%%%%%%%%%%%%%%%%%%%%%%%%%%%%%%%%%%%%%%
%%%%%%%%%%%%%%%%%%%%%%%%%%%%%%%%%کدهای زیر را تغییر ندهید.
\newpage\clearpage

\pagestyle{style2}

\vspace*{-1cm}
\section*{\centering چکیده}
%\addcontentsline{toc}{chapter}{چکیده}
\vspace*{.5cm}
\ffa-abstract
\vspace*{2cm}


{\noindent\large\textbf{واژه‌های کلیدی:}}\par
\vspace*{.5cm}
\fkeywords
%دستور زیر برای شماره گذاری صفحات قبل از فصل اول با حروف ابجد است.
\pagenumbering{alph}
%-----------------------------------------------------------------------------
%در فایل زیر دستورات مربوط به نمایش صفحات فهرست مطالب- فهرست اشکال و جداول است.
%{\pagestyle{style2}
%\tableofcontents}\newpage
%
%\listoffigures
\cleardoublepage
\pagestyle{style6}
\tableofcontents
\pagestyle{style6}
\cleardoublepage
%اگر لیست تصاویر و لیست جداول ندارید ، کدهای زیر را با گذاشتن % در ابتدای آنها، غیرفعال کنید.
\BeginNoToc
\addtocontents{lof}{\lofheading}% add heading to the first page in LoF
\pagestyle{style5}
\listoffigures
\thispagestyle{style5}
\cleardoublepage
\addtocontents{lot}{\lotheading}% add heading to the first page in LoT
\thispagestyle{style4}
\listoftables
\thispagestyle{style4}
%\cleardoublepage
%
\cleardoublepage
\setcounter{savepage}{\arabic{page}}
\mainmatter
\addtocontents{toc}{\tocheading}% add heading to the first page in ToC, after frontmatter entries
\EndNoToc
% در صورت تمایل می‌توانید با فعال کردن دستور بالا، لیست تصاویر را به  پایان‌نامه خود اضافه کنید.
%-------------------------------------------------------------------------symbols(فهرست نمادها)
% وجود لیست نمادها الزامیست.(لطفاً نمادهای خود را جایگذین نمادهای پیش‌فرض کنید.)
%%%%%%%%%%%%%

{\centering\LARGE\textbf{فهرست نمادها}\par}%

\pagenumbering{alph}
\setcounter{page}{\thesavepage}
%\setcounter{page}{6}
\vspace*{1cm}

\pagestyle{style3}
%\thispagestyle{empty}
%\addcontentsline{toc}{chapter}{فهرست نمادها}
\symb{\text{ نماد}}{مفهوم}
\\
%مقادیر بالا را تغییر ندهید
%%%%%%%%%%%%%%%%%%%%%%%%%%%%%%%%%%%%%%%%%%%%%%%%%%%%%%%%%
\symb{V^{SFC}_{i, F}}{
    \small
    مجموعه گره‌های گراف زنجیره‌ی \(i\)ام
}
\symb{E^{SFC}_i}{
    \small
    مجموعه یال‌های گراف زنجیره‌ی \(i\)ام
}
\symb{V_S^{PN}}{
    \small
    مجموعه گره‌های گراف زیرساخت
}
\symb{E_S^{PN}}{
    \small
    مجموعه یال‌های گراف زیرساخت
}
\symb{x_h}{
    \small
    متغیر باینری که نشان می‌دهد زنجیره‌ی \(h\)ام
    پذیرفته شده است یا خیر
}
\symb{y_{wk}}{
    \small
    تعداد نمونه‌هایی از نوع \(k\)
    که روی سرور فیزیکی \(w\) فعال شده‌اند
}
\symb{z^k_{vw}}{
    \small
    متغیر باینری که نشان می‌دهد نمونه‌ی \(v\)
    از نوع \(k\)
    روی سرور فیزیکی \(w\)
    جایگذاری شده است یا خیر
}
\symb{\bar{y}_w}{
    \small
    تعداد نمونه‌هایی از \lr{VNFM} که روی سرور \(w\) فعال شده‌اند
}
\symb{\bar{z}_{hw}}{
    \footnotesize
    متغیر باینری که نشان می‌دهد زنجیره‌ی \(h\) توسط \lr{VNFM}ای که روی سرور \(w\) قرار گرفته است مدیریت می‌شود یا خیر
}
\symb{\tau^{(u,v)}_{ij}}{
    \scriptsize
    متغیر باینری که نشان می‌دهد یال مجازی بین نمونه‌های \(u\) و \(v\) برای نگاشت از یال فیزیکی بین گره‌های \(i\) و \(j\)
    استفاده می‌کند یا خیر
}
\symb{\bar{\tau}^{v}_{ij}}{
    \scriptsize
    متغیر باینری که نشان می‌دهد برای نگاشت ارتباط مدیریتی نمونه‌ی \(v\) از یال فیزیکی بین گره‌های \(i\) و \(j\)
    استفاده شده است یا خیر
}

%%%%%%%%%%%%%%%%%%%%%%%%%%%%%%%%%%%%%%%

\thispagestyle{style3}
\newpage
%\pagestyle{style1}
%%%%%%%%%%%%%%%%%%%%%%%%%%%%%%%%%%%%


\pagenumbering{arabic}
\pagestyle{style1}
%--------------------------------------------------------------------------chapters(فصل ها)
\chapter{مقدمه}


بیشتر سرویس‌های شبکه بر روی سخت افزارهای اختصاصی به نام
\lr{middle box}
ساخته می‌شوند.
تنوع و تعداد رو به افزایش سرویس‌های جدیدی که توسط کاربران تقاضا می‌گردد
باعث هزینه‌های زیاد برای خرید و نگهداری
\lr{middle box}‌ها
توسط اپراتورها شده است.
به تازگی فراهم آورندگان شبکه
شروع به حرکت به سوی مجازی‌سازی و نرم‌افزاری کردن بسترهای شکبه کرده‌اند،
به این ترتیب آن‌ها قادر خواهند بود
سرویس‌های نوآورانه‌ای به کاربران ارائه بدهند.

مجازی‌سازی توابع شبکه‌ راهکاری است که برای همین منظور پیشنهاد شده است.
مجازی‌سازی توابع شبکه‌ در واقع راه‌حل‌های مشخصی را برای چالش‌های جای‌گذاری،
زنجیره‌سازی و هماهنگی سرویس‌های شبکه فراهم می‌آورد.

ایده‌ی اصلی مجازی‌سازی توابع شبکه جداسازی تجهیزات فیزیکی شبکه از کارکردهایی می‌باشد که
بر روی آن‌ها اجرا می‌شوند.
به این معنی که یک کاکرد شبکه مانند دیوار آتش می‌تواند بر روی یک
\lr{TSP}
به عنوان یک نرم‌افزار ساده فرستاده شود.
با این روش یک سرویس می‌تواند به مجموعه‌ای از کارکردهای مجازی شبکه‌ای که می‌توانند به صورت نرم‌افزاری پیاده‌سازی شده
و روی یک یا تعداد سرور استاندارد فیزیکی اجرا شوند، شکسته شود.
کارکرهای مجازی شبکه‌ای می‌توانند در مکان‌های مختلف بازمکان‌یابی یا نمونه‌سازی شوند بدون آنکه
نیاز به خریداری و نصب تجهیز جدیدی باشد.
\cite{Mijumbi2016}

\section{معماری \lr{NFV}}
با توجه به استاندارد \lr{ETSI} معماری \lr{NFV}
از سه عنصر کلیدی تشکیل شده است.

%--------------------------------------------------------------------------appendix( مراجع و پیوست ها)
\chapterfont{\vspace*{-2em}\centering\LARGE}%

\appendix
\bibliographystyle{plain-fa}
\bibliography{references}
% \chapter*{‌پیوست}
\markboth{پیوست}{}
\addcontentsline{toc}{chapter}{پیوست}
موضوعات مرتبط با متن گزارش پایان نامه كه در يكی از گروه‌های زير قرار می‌گيرد، در بخش پيوست‌ها آورده شوند:
\begin{enumerate}
\item  اثبات های رياضی يا عمليات رياضی طولانی‌.‌
\item داده و اطلاعات نمونه (های) مورد مطالعه (\lr{Case Study}) چنانچه طولانی باشد‌.‌
\item نتايج كارهای ديگران چنانچه نياز به تفصيل باشد‌.‌
\item مجموعه تعاريف متغيرها و پارامترها، چنانچه طولانی بوده و در متن به انجام نرسيده باشد‌.‌
\end{enumerate}
% براي شماره‌گذاري روابط، جداول و اشكال موجود در پيوست‌ از ساختار متفاوتي نسبت به متن اصلي استفاده مي‌شود كه در زير به‌عنوان نمونه نمايش داده شده‌است. 
% \begin{equation}
%F=ma
%\end{equation}
\section*{کد میپل }
\begin{latin}
\begin{verbatim}

with(DifferentialGeometry):
with(Tensor):
DGsetup([x, y, z], M)
																	frame name: M
a := evalDG(D_x)
																	D_x
b := evalDG(-2 y z D_x+2 x D_y/z^3-D_z/z^2)


\end{verbatim}
\end{latin}
%--------------------------------------------------------------------------dictionary(واژه نامه ها)
%اگر مایل به داشتن صفحه واژه‌نامه نیستید، خط زیر را غیر فعال کنید.
\parindent=0pt
%
\chapter*{واژه‌نامه‌ی فارسی به انگلیسی}
\pagestyle{style9}

\addcontentsline{toc}{chapter}{واژه‌نامه‌ی فارسی به انگلیسی}
%%%%%%
\begin{multicols*}{2}

{\bf آ}
%%\vspace*{3mm}

\vspace*{3mm}
{\bf ب}
%%\vspace*{3mm}

\vspace*{3mm}
{\bf پ}
%%\vspace*{3mm}

\vspace*{3mm}
{\bf ت}
%%\vspace*{3mm}

\vspace*{3mm}
{\bf ث}
%%\vspace*{3mm}

\vspace*{3mm}
{\bf ج}
%%\vspace*{3mm}

\vspace*{3mm}
{\bf چ}
%%\vspace*{3mm}

\vspace*{3mm}
{\bf ح}
%%\vspace*{3mm}

\vspace*{3mm}
{\bf خ}
%%\vspace*{3mm}

\vspace*{3mm}
{\bf د}
%%\vspace*{3mm}

\vspace*{3mm}
{\bf ر}
%%\vspace*{3mm}

\vspace*{3mm}
{\bf ز}
%%\vspace*{3mm}

\vspace*{3mm}
{\bf س}
%%\vspace*{3mm}

\vspace*{3mm}
{\bf ص}
%%\vspace*{3mm}

\vspace*{3mm}
{\bf ض}
%%\vspace*{3mm}

\vspace*{3mm}
{\bf ط}
%%\vspace*{3mm}

\vspace*{3mm}
{\bf ظ}
%%\vspace*{3mm}

\vspace*{3mm}
{\bf ع}
%%\vspace*{3mm}

\vspace*{3mm}
{\bf ف}
%%\vspace*{3mm}

\farsiTOenglish{فراهم‌آورنده‌ی شبکه}{Network Provider}



\vspace*{3mm}
{\bf ک}
%%\vspace*{3mm}

\vspace*{3mm}
{\bf گ}
%%\vspace*{3mm}

\vspace*{3mm}
{\bf م}
%%\vspace*{3mm}

\farsiTOenglish{مجازی‌سازی کارکردهای شبکه}{Network Function Virtualization}


\vspace*{3mm}
{\bf ن}
%%\vspace*{3mm}

\vspace*{3mm}
{\bf و}
%%\vspace*{3mm}

{\bf ه}
%%\vspace*{3mm}

{\bf ی}
%%\vspace*{3mm}


\end{multicols*}%
% %%%%%%
\chapter*{واژه‌نامه}
\pagestyle{style9}
\lhead{\thepage}\rhead{واژه‌نامه}
\addcontentsline{toc}{chapter}{واژه‌نامه}

\englishTOfarsi{Interface}{رابط}

\englishTOfarsi{Application}{کاربرد}

\englishTOfarsi{Optimality Gap}{شکاف بهینه}

\englishTOfarsi{Edge}{یال}

\englishTOfarsi{Node}{گره}

\englishTOfarsi{Virtualization}{مجازی‌سازی}

\englishTOfarsi{Constraint}{محدودیت}

\englishTOfarsi{Heuristic}{مکاشفه‌ای}

\englishTOfarsi{Offline}{برون‌خط}

\englishTOfarsi{Online}{برخط}

\englishTOfarsi{Image}{تصویر}

\englishTOfarsi{Classifier}{دسته‌بند}

\englishTOfarsi{Forwarding}{جلورانی}

\englishTOfarsi{Function}{کارکرد}

\englishTOfarsi{Pod}{غلاف}

\englishTOfarsi{Aggregation}{تجمعی}

\englishTOfarsi{Edge}{لبه}

\englishTOfarsi{Instance}{نمونه}

\englishTOfarsi{Feasible Set}{مجموعه امکان‌پذیر}
%--------------------------------------------------------------------------index(نمایه)
%اگر مایل به داشتن صفحه نمایه نیستید، خط زیر را غیر فعال کنید.
\pagestyle{style7}
\printindex
\pagestyle{style7}
%کلمات کلیدی انگلیسی
% \latinkeywords{NFV, SFC, Optimization, ILP}
%چکیده انگلیسی

\en-abstract{
In the old times, Network providers use hardware network functions to create their service chains, but a change in this manner is difficult and may cause many service distribution. SFC and NFV is the solution to this difficulty. By using SFC and NFV, providers can provision chains dynamically and then change them in runtime. One of the main requirements is management and monitoring for the chains.
In this research, we consider the chain acceptance problem subject to management resources. In the first step, we formulate
problem with ILP and then implement it in CPLEX framework. As we know, ILP problems are NP-Hard, so we need an NP solution to the problem. In this research, we create a heuristic algorithm and compare its result with the optimal solution. In the end, the heuristic solution produces near-optimal results in the polynomial time.
}
%%%%%%%%%%%%%%%%%%%%% کدهای زیر را تغییر ندهید.

\newpage
\thispagestyle{empty}
\begin{latin}
\section*{\LARGE\centering Abstract}

\een-abstract

\vspace*{.5cm}
{\large\textbf{Key Words:}}\par
\vspace*{.5cm}
\elatinkeywords
\end{latin}

% در این فایل، عنوان پایان‌نامه، مشخصات خود و چکیده پایان‌نامه را به انگلیسی، وارد کنید.
%%%%%%%%%%%%%%%%%%%%%%%%%%%%%%%%%%%%
\baselineskip=.6cm
\begin{latin}

\latinfaculty{Department of Computer Engineering \& Information Technology}


\latintitle{Virtualized Network Service Function chaining Subject to Management Resource Constraint}


\firstlatinsupervisor{Dr.\ Bahador Bakhshi}

%\secondlatinsupervisor{Second Supervisor}

% \firstlatinadvisor{Dr. }

%\secondlatinadvisor{Second Advisor}

\latinname{Parham}

\latinsurname{Alvani}

\latinthesisdate{April 2018}

\latinvtitle{}
\end{latin}
\end{document}